\section{Proof of Lam's determinant conjecture}\label{sec:determinants}

\begin{proof}[Proof of Theorem~\ref{thm:Lam-real}]

Theorem~\ref{thm:deRham-determinant},
Corollary~\ref{cor:opposite-periods}, and
\eqref{eq:det-period-relation} give
\[
 \det B_b^{\cB}
 =\pm(2\pi i)^{-r\beta(M)}
       R_M(\mathbf a)\Betafunc_M(\mathbf a)\Betafunc_M(-\mathbf a).
\]
Substituting the products defining $R_M$ and $\Betafunc_M$ gives
\[
 \det B_b^{\cB}
 =\pm(2\pi i)^{-r\beta(M)}
 \frac{\displaystyle
   \prod_{F\in\cL\setminus(\cL_0\cup\{\varnothing\})}
      \bigl(a_F^{-1}\Gamma(1+a_F)\Gamma(1-a_F)\bigr)^{w_F}}
      {\displaystyle
   \prod_{F\in\cL_0\setminus\{E\}}
      \bigl(a_F^{-1}\Gamma(1+a_F)\Gamma(1-a_F)\bigr)^{w_F}}.
\]
The reflection formula and $b_F=e^{-\pi i a_F}$ give
\begin{equation}\label{eq:reflection}
 a_F^{-1}\Gamma(1+a_F)\Gamma(1-a_F)
 =\frac{\pi}{\sin(\pi a_F)}
 =\frac{2\pi i b_F}{1-b_F^2}.
\end{equation}
Collecting the powers of $2\pi i$ gives
\[
 \det B_b^{\cB}
 =\pm(2\pi i)^\delta\frac{\displaystyle
      \prod_{F\in\cL\setminus(\cL_0\cup\{\varnothing\})}
          \bigl(b_F/(1-b_F^2)\bigr)^{w_F}}
        {\displaystyle
      \prod_{F\in\cL_0\setminus\{E\}}
          \bigl(b_F/(1-b_F^2)\bigr)^{w_F}}.
\]
Here $\delta$ is the total exponent of $2\pi i$.
For $F\in\cL_0$, the relation $b_E=1$ gives
\[
 1-b_F^2=-b_F^2(1-b_{E\setminus F}^2).
\]
Using Proposition~\ref{prop:Betti-Varchenko} and absorbing the signs
into $c=\pm(2\pi i)^\delta$ gives
\[
 \det\bigl((V_b^{-1})^{\cB}\bigr)
 =c\left(\prod_{\varnothing\ne F\ne E}b_F^{w_F}\right)
 \frac{\displaystyle\prod_{F\in\cL_0\setminus\{E\}}
            (1-b_{E\setminus F}^2)^{w_F}}
      {\displaystyle\prod_{F\in\cL\setminus(\cL_0\cup\{\varnothing\})}
            (1-b_F^2)^{w_F}}.
\]
By \cite[Lemma 2.15]{LopezdeMedranoRinconShaw2020}, the sum
\[
 s=\sum_{\substack{F\in\cL\setminus\{\varnothing,E\}\\i\in F}}w_F
\]
is independent of $i\in E$. This also follows from
\cite[Corollary 2.1.15]{Gioan2002}. Hence the monomial factor is
\[
 \prod_{\varnothing\ne F\ne E}b_F^{w_F}
 =\prod_{i\in E}b_i^s=b_E^s=1.
\]
This proves \eqref{eq:Lam-determinant} up to the constant $c$. After eliminating
$b_0$, both sides are rational functions of the finite parameters
$b_i$, $i\in E_0$. Their equality on nonresonant weights therefore
gives an identity of rational functions. Both sides are regular at
$b_i=0$ for all $i\in E_0$, where $V_b$ becomes the identity and the
product becomes $1$. Thus $c=1$, proving the theorem.
\end{proof}

We obtain the following identity as a corollary.
\begin{corollary}\label{cor:degree}
For an affine oriented matroid $(\M,0)$ of rank $r+1$,
\[
 \sum_{F\in\cL\setminus(\cL_0\cup\{\varnothing\})}w_F
 -\sum_{F\in\cL_0\setminus\{E\}}w_F=r\beta(M).
\]
\end{corollary}
\begin{proof}
For real arrangements, $\delta$ in the preceding proof is the
left-hand side minus $r\beta(M)$. Since $1=|c|=(2\pi)^\delta$,
we have $\delta=0$.
For arbitrary affine oriented matroids,
\eqref{eq:combinatorial-deRham} and \eqref{eq:nbc-form} show that
each entry of $D_{\mathbf a}^{\mathrm{nbc}}$ is homogeneous of degree
$r$ in the weights. Its determinant has degree $r\beta(M)$. Comparing with
$R_M(\mathbf a)^{-1}$ in Theorem~\ref{thm:deRham-determinant}
gives the same identity.
\end{proof}

\begin{example}\label{ex:line4}
We conclude by computing the period relation for the two-point
arrangement. Substituting Examples~\ref{ex:line2} and~\ref{ex:line3}
into \eqref{eq:det-period-relation} and applying~(\ref{eq:reflection}) gives
\[
 \begin{aligned}
 \det B_b^{\cB} = B_b(P,P)
 &=\frac{a_0}{2\pi i a_1a_2}
   \frac{\Gamma(1+a_1)\Gamma(1-a_1)
         \Gamma(1+a_2)\Gamma(1-a_2)}
        {\Gamma(1+a_0)\Gamma(1-a_0)}\\
 &=\frac{b_1 b_2(1-b_0^2)}{b_0 (1-b_1^2)(1-b_2^2)}.
 \end{aligned}
\]
Since $b_0b_1b_2=1$,
\[
 \begin{aligned}
 \det B_b^{\cB}
 &=\frac{(1-b_0^2)}{b_0^2 (1-b_1^2)(1-b_2^2)}\\
 &=-\frac{1-b_1^2b_2^2}{(1-b_1^2)(1-b_2^2)}
 =-(V_b^{-1})^{\cB},
 \end{aligned}
\]
and we recover the Betti entry in Example~\ref{ex:line}.
\end{example}
